\section{Conclusions}\label{sec:conclusions}

The angle bias of single-time-per-channel track reconstruction, reported by every strip micro-TPC that compared itself with
a set angle and attributed to capacitive coupling, is a geometric property of reading a short drift column one time per
channel. A time per channel is a functional of the charge of every depth that reaches the channel; the channels at the two
ends of the column receive charge of the wrong depth from the partial overlap of the track with the strip and from any
lateral spread, and rotate the fitted line toward a steeper track. The ratio of the reconstructed to the true tangent is
$1 + C/N^2 + b$, with $N$ the strips in the column, $C$ set by the lateral spread in units of the pitch and the estimator, and
$b$ an interior compression from the growth of the spread with depth that is small for a leading edge and absent for a
charge-weighted time.

On a 28~mm resistive-strip Micromegas with a per-track silicon reference, $C = 7.0$ (X) and $11.7$ (Y) with $|b| < 0.02$; a
forward model of the waveforms reproduces $C$ to 2\,\% and decomposes it into strip geometry (2.3), the same-time footprint
(+1.5), drift diffusion (+3.1) and, across the resistive strips, RC spreading (+4.6); the published strip-$\mu$TPC points fall
between the geometry-only and the full curves in the order of their lateral spread; three further chambers, whose resistive
spreading constants span a factor three, reproduce the model's $C$ within 10\,\% on both views and order with the spreading
constant on the resistive view. Depth-resolved stacking of the strip
waveforms separates the same-time share at the mesh from a growth that matches a transverse diffusion of
\SI{350}{\micro\metre\per\sqrt{cm}}, and finds a same-time second-neighbour share of 2.4\,\% per side, mostly electronic, that no
Gaussian footprint gives. The position-in-strip test bounds a fixed nearest-neighbour coupling to about 6\,\% on the view
across the resistive strips, a factor three to five below the published fractions, and its closure shows the fitted position
to be adequate for the test. Nothing in the data requires a coupling of the published size, and nothing in the angle could
have told the difference.

The literature's remedies address the ends. Scored on the same tracks, the edge cuts and shifts remove a quarter of the
bias, the neighbour subtraction at its tuned factor of 30\,\% removes two thirds on the metallic-like view and harms the
resistive one, and the charge-weighted time helps only after a subtraction; the best hit-level recipe leaves the angle
resolution at 1.5 to 2 times that of a per-event forward fit of the waveforms, which is unbiased to 1\,\% at every angle
because the spreads are in its model. The tuned subtraction factor is a measure of the total same-time spread, not of a
capacitance.

For a detector builder the useful numbers are two: the strips in the column and the lateral spread in pitches. A column of
more than about 25 strips, or a charge-weighted time on a column with no ghosts, hides the effect; a column of a few strips
read by one time each will show it at the 10 to 50\,\% level, whatever the front end, and must either be calibrated against
an external angle at every inclination or be read as waveforms and fitted. A $\mu$TPC cannot measure its own drift velocity
from the ladder, and a coupling fraction tuned on the angle cannot be quoted as a property of the board.

\section*{Acknowledgements}
\tbw{MX17 collaboration, n\_TOF, UCLA group; the literature survey drew on the open theses of the LMU and NTUA groups.}
